\documentclass[11pt,a4paper]{article}
\usepackage[T1]{fontenc}
\usepackage{lmodern,microtype}
\usepackage[margin=25mm]{geometry}
\usepackage{amsmath,amssymb,amsthm,mathtools}
\usepackage{booktabs,array,longtable,enumitem,xcolor,fancyhdr,listings}
\usepackage[colorlinks=true,linkcolor=blue!45!black,citecolor=blue!45!black,urlcolor=blue!45!black,hypertexnames=false]{hyperref}
\setlength{\emergencystretch}{3em}
\setlength{\parskip}{.25em}
\setlength{\headheight}{14pt}
\pagestyle{fancy}\fancyhf{}
\fancyhead[L]{\small Dyadic saturation of vanishing sign sums}
\fancyhead[R]{\small Research manuscript}\fancyfoot[C]{\thepage}
\newtheorem{theorem}{Theorem}[section]
\newtheorem{lemma}[theorem]{Lemma}
\newtheorem{proposition}[theorem]{Proposition}
\newtheorem{corollary}[theorem]{Corollary}
\theoremstyle{definition}\newtheorem{definition}[theorem]{Definition}
\theoremstyle{remark}\newtheorem{remark}[theorem]{Remark}
\newcommand{\N}{\mathbb N}\newcommand{\Z}{\mathbb Z}\newcommand{\R}{\mathbb R}
\newcommand{\Dset}{\mathcal D}\newcommand{\wt}{\operatorname{wt}}
\newcommand{\LL}{\mathcal L}\newcommand{\UU}{\mathcal U}
\DeclareMathOperator{\ord}{ord}
\lstset{basicstyle=\ttfamily\scriptsize,breaklines=true,columns=fullflexible,keepspaces=true,showstringspaces=false,frame=single,rulecolor=\color{black!20},aboveskip=8pt,belowskip=8pt}
\title{Dyadic Saturation of Vanishing Periodic Sign Sums\\[5pt]
\large Explicit cycle constructions, an exact period bound,\\and companion graph and spectral results}
\author{Keio Research Notes\\\small Developed with computational and independent proof checks}
\date{2 November 2026}
\hypersetup{pdftitle={Dyadic Saturation of Vanishing Periodic Sign Sums},pdfauthor={Keio Research Notes}}
\begin{document}
\maketitle
\begin{abstract}
Let $\Dset$ be the positive integers $N$ admitting an $N$-periodic function
$f:\Z\to\{-1,1\}$ with zero mean and $\sum_{n\ge1}f(n)/n=0$.
For every odd integer $m\ge3$, we construct such a function with
$N=2^a m$ and $a\le5\lceil\log_2m\rceil+10$; in particular
$N\le32768m^6$. A cyclic-word construction balances the exact
$2$-adic valuation weights of $H(mu)-mH(u)$. Signed binary expansions
transfer the balance to periodic signs, and an elementary Dirichlet-series
identity proves exact vanishing. No numerical approximation to the
infinite sum is used in the existence proof. Together with the exclusion
of powers of two, this classifies the dyadic saturation of $\Dset$ and
settles the prime-lift conjecture explicitly left open in corpus item
[093]. It does not classify $\Dset$ itself. Independent appendices
prove a factorial order bound for supportive decorated graph cores and
give an exact annular counterexample to the common-smallness question
for first Steklov eigenfunctions as printed in a 2020 problem list,
together with a sharp normalized replacement for multiplicity two.
A further appendix constructs exact randomized-time uniform mixing
on every path $P_{p^k-1}$, where $p$ is an odd prime and $k\ge2$.
The manuscript distinguishes these mathematical conclusions from
unverified claims of historical priority or comparative research value.
\end{abstract}

\section{Results, provenance, and scope}
For a periodic function of mean zero, write
\[
 \LL(f)=\sum_{n=1}^{\infty}\frac{f(n)}n.
\]
The series converges by summation by parts. Define
\begin{equation}\label{eq:Ddef}
 \Dset=\left\{N\ge1:\ \exists f:\Z\to\{-1,1\},\quad
 f(n+N)=f(n),\quad \sum_{n=1}^N f(n)=0,\quad \LL(f)=0\right\}.
\end{equation}
Here $N$ is a period, not necessarily the least period. This convention
matters: $\Dset$ is closed under taking positive multiples.

Murty and Pathak ask for a classification of the periods in
\eqref{eq:Ddef}; their Section 4.4 records Tengely's example of period
$36$ and rules out $N\equiv2\pmod4$ \cite{MP2022}. This is different
from the classical Erd\H{o}s problem in which the function is zero,
rather than a sign, on multiples of the period. We retain that distinction
throughout.

The immediate starting point is corpus item [093] \cite{Corpus093}.
Its Theorem E supplies a signed-digit construction, its Theorem C treats
Fermat primes, and its Conjecture 9.1 asks whether
\[
 a_0(p):=\min\{a\ge0:2^ap\in\Dset\}<\infty
 \quad\hbox{for every odd prime }p.
\]
The same item reports an unsuccessful search at $p=127$. Our main result
removes that existence obstruction for every odd multiplier, including
composite multipliers.

\begin{theorem}[Explicit dyadic lifts]\label{thm:main}
For every odd integer $m\ge3$, there are an integer $a$ and an explicitly
constructible $2^am$-periodic sign function $f$ such that
\begin{equation}\label{eq:main}
 a\le5\lceil\log_2m\rceil+10,\qquad
 \sum_{n=1}^{2^am}f(n)=0,\qquad \LL(f)=0.
\end{equation}
The construction can therefore be made with $2^am\le32768m^6$.
\end{theorem}

\begin{corollary}[Exact dyadic saturation]\label{cor:saturation}
For a positive integer $n$, some $2^jn$, $j\ge0$, belongs to $\Dset$
if and only if $n$ is not a power of two. In particular, $a_0(p)$ is
finite for every odd prime $p$, and
\[
 \{a\ge0:2^ap\in\Dset\}=\{a\ge a_0(p)\}.
\]
\end{corollary}

The existence part of Theorem~\ref{thm:main} is elementary and
self-contained. Only the exclusion of pure powers of two in
Corollary~\ref{cor:saturation} uses a published necessary condition of
Okada, in the form stated by Murty and Pathak \cite[Theorem 1.1]{MP2022}.
The proof is supplied in Section~\ref{sec:consequences}.

\paragraph{What changes relative to the source.}
The new step relative to [093] is an exact cyclic-word filling argument.
It replaces a search for a suitable digit tree by a uniform construction
on sufficiently long multiplication cycles. The transfer to signs is
reproved directly, and it works for every odd integer $m$, so no
assumption of primality is hidden in the proof. The exponent bound is
deliberately coarse; it is an existence bound, not a formula for $a_0$.

\paragraph{Snapshot and verification status.}
The inspected repository contains numbered items [001]--[102]. Its
commit is
\begin{center}\small\texttt{ecd87de3348dcb706a76bb12fde722d7a348a012}.\end{center}
The commit timestamp is 2 November 2026, 19:38:17, UTC+09:00.
The repository describes its contents as research drafts rather than
peer-reviewed or formally verified theorems. We do not infer correctness
from a title or from inclusion in that repository. The existence theorem
below includes its required proofs. The companion graph theorem is also
proved independently; its subsequent radius application explicitly
identifies the earlier reduction it uses.

\section{The finite cyclic-word problem}
For a nonzero integer $z$ put
\[
 \wt(z)=2^{v_2(z)}.
\]
We first solve a finite problem: construct a cyclic word of odd integers
$h_0,\ldots,h_{t-1}$ for which every $h_{i+1}-mh_i$ is nonzero and
\begin{equation}\label{eq:wordbudget}
 \sum_{i=0}^{t-1}\wt(h_{i+1}-mh_i)
       =(m-1)\sum_{i=0}^{t-1}\wt(h_i).
\end{equation}
Indices are read modulo $t$. Since all $h_i$ will be odd, the right-hand
side is simply $(m-1)t$.

\begin{lemma}[An exact cyclic budget]\label{lem:word}
Let $m\ge3$ be odd, and set
\begin{equation}\label{eq:parameters}
 \ell=\lceil\log_2m\rceil,\quad r=v_2(m-1),\quad
 s=v_2(m^2-1),\quad t=2^{\ell+3},\quad b=2^r.
\end{equation}
There is a cyclic word satisfying \eqref{eq:wordbudget}, all of whose
entries belong to $\{1,-1\}\cup\{m+2^M:M>s\}$. It can be chosen with
\begin{equation}\label{eq:wordbounds}
 \max_i|h_i|<2^{2\ell+4},\qquad
 \max_i|h_{i+1}-mh_i|<2^{3\ell+5}.
\end{equation}
\end{lemma}
\begin{proof}
Start with the constant word $h_i=1$, whose total weight is $tb$.
At an isolated odd position replace $1$ by $A$, leaving its neighbors
equal to $1$. This changes the total weight by
\begin{equation}\label{eq:excursion}
 \wt(A-m)+\wt(1-mA)-2b.
\end{equation}
If $A=m+2^M$ with $M>s$, the two weights are $2^M$ and $2^s$.
Indeed $1-mA=1-m^2-m2^M$ has valuation $s$ because the two summands
have different valuations. Thus \eqref{eq:excursion} equals
$2^M+C$, where $C=2^s-2b$.

Suppose first that $m\equiv1\pmod4$. Then $s=r+1$, so $C=0$.
The missing weight is
\[
 D=t(m-1-b)\ge0.
\]
If $D=0$, retain the constant word. Otherwise, for each set bit $2^M$
in the binary expansion of $D$, insert one excursion with $A=m+2^M$.
Every such exponent is greater than $s$: $D$ is divisible by
$2^{\ell+3+r}$ and $\ell+3+r>r+1=s$. Also
$D<2^{2\ell+3}$, so there are at most $2\ell+3$ excursions. For
$m\equiv1\pmod4$, $m\ge5$ and $\ell\ge3$; hence
$2(2\ell+3)\le t$. Place these entries at successive odd positions.
Their incident edge pairs are disjoint, and their increments sum to $D$.

For $m=3$ the constant word already works, because $b=m-1=2$.
It remains to treat $m\equiv3\pmod4$, $m\ge7$. Here $b=2$ and
$C=2^s-4>0$. An excursion with $A=-1$ has both edge weights
$2^{v_2(m+1)}=2^{s-1}$, and therefore contributes exactly $C$.
Put
\[
 k=2^{\ell+1},\qquad D=t(m-3),\qquad B=D-kC.
\]
Since $2^{s-1}\mid m+1\le2^\ell$, we have $s\le\ell+1$ and
$C\le2m-2$. Consequently
\[
 B\ge2^{\ell+1}\{4(m-3)-(2m-2)\}
   =2^{\ell+2}(m-5)>0.
\]
Both $D$ and $kC$ are divisible by $2^{s+1}$: $D$ contains
$2^{\ell+3}$, while $C=4(2^{s-2}-1)$ contains exactly $2^2$.
Thus every bit in $B$ has exponent $M>s$. Moreover
$B<D<2^{2\ell+3}$, so $B$ has at most $2\ell+3\le k$ set bits.
Use one high excursion $A=m+2^M$ for each such bit and use $A=-1$
for the remaining excursions, making exactly $k$ excursions.
Their total increment is $B+kC=D$, as required. They fit because
$2k=t/2\le t$.

No edge difference in any of these words vanishes: its weight has
just been computed, and the unchanged differences are $1-m\ne0$.
Every binary power used is less than $tm<2^{2\ell+3}$; hence
$|h_i|<2^{2\ell+4}$. Since $m+1\le2^\ell$, the elementary estimate
$|h_{i+1}-mh_i|\le(m+1)\max|h_i|$ gives the looser second bound
in \eqref{eq:wordbounds}.
\end{proof}

\section{Multiplication cycles and congruence compatibility}
Write $\UU_j=(\Z/2^j\Z)^\times$ for $j\ge1$ and let $\UU_0$
be a one-point set. The following step turns the word into a function
on units, with the congruence conditions needed for signed digits.

\begin{lemma}[Paired cycles]\label{lem:cycles}
With the parameters in \eqref{eq:parameters}, put
\[
 J=\ell+s+2,\qquad q=2^J.
\]
Multiplication by $m$ on $\UU_J$ has all cycles of length $t$, and
negation pairs distinct cycles. There is a function $H:\UU_J\to\Z$
such that, on extending it periodically to odd integers and putting
$G(u)=H(mu)-mH(u)$, the following hold:
\begin{align}
 &H(-u)=-H(u),\quad H(u)\ne0,\quad G(u)\ne0,\label{eq:anti}\\
 &H(u)\equiv H(v)\pmod{2^j}\quad\hbox{if }u\equiv v\pmod{2^j},\label{eq:lipschitz}\\
 &\sum_{u\in\UU_J}\wt(G(u))
       =(m-1)\sum_{u\in\UU_J}\wt(H(u)).\label{eq:unitbudget}
\end{align}
The congruence assertion holds for every $j\ge1$.
\end{lemma}
\begin{proof}
For even $e$, the elementary lifting identity gives
$v_2(m^e-1)=s+v_2(e)-1$. It follows that
$\ord_{2^J}(m)=2^{J-s+1}=2^{\ell+3}=t$. This identity can also be
obtained by repeatedly factoring $m^{2e}-1$; for odd $e$ the valuation
of $m^e-1$ is $v_2(m-1)$.

If $m\equiv1\pmod4$, no power is $-1\pmod{2^J}$.
If $m\equiv3\pmod4$, an even power is $1\pmod4$, and an odd power
$m^e$ satisfies $v_2(m^e+1)=v_2(m+1)=s-1<J$. Thus $-1$ never lies
in the subgroup generated by $m$. All cycles are consequently paired
with distinct negative cycles.

Choose one cycle in each pair, choose its initial unit $u$, and use the
word of Lemma~\ref{lem:word} to set
\[
 H(m^iu)=q h_i,\qquad H(-m^iu)=-q h_i\quad(0\le i<t).
\]
This is well-defined. Every value is a multiple of $q$, so
\eqref{eq:lipschitz} holds when $j\le J$; for $j>J$, the two values
are identical. On each positive cycle, $G(m^iu)=q(h_{i+1}-mh_i)$,
and the negative cycle has the opposite values. The word budget,
multiplied by $q$ and summed over the pairs, proves
\eqref{eq:unitbudget}. All nonvanishing assertions follow from the word.
\end{proof}

Lifting $H,G$ to $\UU_a$, $a\ge J$, multiplies both sums in
\eqref{eq:unitbudget} by $2^{a-J}$. By \eqref{eq:wordbounds},
\[
 M:=\max_{u}\{|H(u)|,|G(u)|\}<2^{J+3\ell+5}.
\]
The choice
\begin{equation}\label{eq:adepth}
 a=J+3\ell+7
\end{equation}
therefore ensures
\begin{equation}\label{eq:size}
 4M+1<2^{a+1}.
\end{equation}
Because $s\le\ell+1$, this choice satisfies $a\le5\ell+10$.
One may use any smaller $a\ge J$ satisfying \eqref{eq:size}; this
is how the numerical examples below are shortened.

\section{Signed binary trees and their exact layer sums}
Every odd integer $x$ with $|x|<2^{a+1}$ has a unique expansion
\begin{equation}\label{eq:signed}
 x=\sum_{i=0}^{a}2^i d_i(x),\qquad d_i(x)\in\{-1,1\}.
\end{equation}
Indeed, $(x+2^{a+1}-1)/2$ is an integer between $0$ and
$2^{a+1}-1$; its binary digits $b_i$ give $d_i=2b_i-1$.
Digits $0,\ldots,j$ agree for two inputs if those inputs agree
modulo $2^{j+2}$.

\begin{definition}
A signed digit tree of depth $a$ is a family of functions
$d_i:\UU_i\to\{-1,1\}$, $0\le i\le a$. Its value function is
\[
 X(u)=\sum_{i=0}^a2^i d_i(u\bmod2^i),\qquad u\in\UU_a.
\]
Its associated $2^a$-periodic sign function $A$ is defined by
\[
 A(2^{a-i}u)=d_i(u),\quad u\in\UU_i,
\]
where $i=0$ denotes the zero residue modulo $2^a$.
\end{definition}

Suppose $H$ satisfies \eqref{eq:lipschitz} and \eqref{eq:size}.
For $\kappa\in\{-1,1\}$ the numbers
$X(u)=\kappa+4H(u)$ have consistent signed expansions: the digit $d_i$
depends only on $u\bmod2^i$, and the digits $d_0,d_1$ are constant.
For $i\ge1$ this follows from congruence modulo $2^{i+2}$, and
for $i=0$ it follows from the common residue modulo $4$.
In particular $d_0=-\kappa$.

For any tree, set
\[
 u_0=u_1=1,\quad u_i=2^{i-1}\ (i\ge1),\qquad
 S_i=\sum_{v\in\UU_i}d_i(v),\qquad
 \overline X=\frac1{2^{a-1}}\sum_{u\in\UU_a}X(u).
\]
Counting the lifts from $\UU_i$ to $\UU_a$ gives
\begin{equation}\label{eq:mean}
 \overline X=d_0+2\sum_{i=1}^a S_i,
 \qquad \sum_{n=1}^{2^a}A(n)=\sum_{i=0}^aS_i
       =\frac{\overline X+d_0}{2}.
\end{equation}
Thus an antisymmetric $H$ and $X=\kappa+4H$ always produce a
mean-zero sign function, since $\overline X=\kappa$ and $d_0=-\kappa$.

The logarithmic layer statistic is
\begin{equation}\label{eq:E}
 E(A)=aS_0+\sum_{i=1}^a(a-i-1)S_i.
\end{equation}
We next compute it directly from the valuation weights of $H$.

\begin{lemma}[Signed-digit pairing]\label{lem:pairing}
Define
\[
 \lambda_a(x)=a2^{1-a}d_0(x)
       +\sum_{i=1}^a(a-i-1)2^{i-a}d_i(x).
\]
If $z\ne0$, $\kappa=\pm1$, and $4|z|+1<2^{a+1}$, then
\begin{equation}\label{eq:pairing}
 \lambda_a(\kappa+4z)+\lambda_a(\kappa-4z)
       =-\kappa\,2^{v_2(z)+3-a}.
\end{equation}
Consequently, for $X=\kappa+4H$ with antisymmetric nonzero $H$,
\begin{equation}\label{eq:Eweight}
 E(A)=-\kappa\,2^{2-a}\sum_{u\in\UU_a}\wt(H(u)).
\end{equation}
\end{lemma}
\begin{proof}
Let $x_k=\sum_{i<k}2^i d_i(x)$. Direct summation gives
\[
 2^a\lambda_a(x)=\sum_{k=1}^a x_k-x+(a+1)d_0(x).
\]
Write $v=v_2(z)$, $x=\kappa+4z$, $y=\kappa-4z$.
Then $d_0(x)=d_0(y)=-\kappa$ and $x_1=y_1=-\kappa$.
For $k\ge2$, the unique representatives $x_k,y_k$ in
$(-2^k,2^k)$ have sum $2\kappa-\kappa2^{k+1}$ when $k\le v+1$,
and sum $2\kappa$ when $k>v+1$. This follows by reducing $x,y$
modulo $2^{k+1}$ in the signed expansion: the exceptional case is
exactly $x\equiv y\equiv\kappa\pmod{2^{k+1}}$.
The size assumption implies $v\le a-2$. Summing these expressions
therefore yields \eqref{eq:pairing}.

Counting lifts in each digit shows
$E(A)=\sum_{u\in\UU_a}\lambda_a(X(u))$. Pairing $u$ with $-u$
and using \eqref{eq:pairing} proves \eqref{eq:Eweight}; each pair
has the same valuation twice, which accounts for the factor $2^{2-a}$.
\end{proof}

\section{An elementary cancellation identity}\label{sec:transfer}
The next lemma applies to arbitrary real layer coefficients, not only
signs. It is the analytic transfer needed below.

\begin{lemma}[Constant value function]\label{lem:cancellation}
Let $a\ge1$ and let $C$ be a mean-zero, $2^a$-periodic function. Define its layer
functions $c_i$ and sums $S_i$ as above, and suppose
\[
 X_C(u)=\sum_{i=0}^a2^i c_i(u\bmod2^i)
\]
is constant on $\UU_a$. Then
\begin{equation}\label{eq:cancellation}
 \LL(C)=-2^{-a}E(C)\log2.
\end{equation}
\end{lemma}
\begin{proof}
Equation \eqref{eq:mean}, which is linear and remains valid for real
coefficients, shows that the constant is $-c_0$. For real $z>1$ set
\[
 D_i(z)=\sum_{\substack{n\ge1\\n\text{ odd}}}\frac{c_i(n)}{n^z}
 \quad(i\ge1).
\]
Grouping the series for $C$ by its dyadic valuation gives
\[
 2^{az}\sum_{n\ge1}\frac{C(n)}{n^z}
 =c_0\zeta(z)+\sum_{i=1}^a2^{iz}D_i(z).
\]
Since $\sum_{i\ge1}2^i c_i(n)=-2c_0$ for odd $n$, the right side
equals
\begin{equation}\label{eq:abel}
 c_0(2^{1-z}-1)\zeta(z)
   +\sum_{i=1}^a(2^{iz}-2^i)D_i(z).
\end{equation}
Periodicity gives $\lim_{z\downarrow1}(z-1)D_i(z)=S_i/2^i$.
For example, subtract the mean $S_i/2^i$ from the $2^i$-periodic
coefficient sequence supported on odd integers; the remainder has
bounded partial sums. Also $(z-1)\zeta(z)\to1$. Taking the limit in
\eqref{eq:abel} yields
\[
 2^a\LL(C)=\left(-c_0+\sum_{i=1}^a iS_i\right)\log2.
\]
The same bounded-partial-sum argument justifies the limit of the
mean-zero series on the left. Since $\sum_iS_i=0$,
$E(C)=c_0-\sum_{i\ge1}iS_i$, proving the assertion.
\end{proof}

\begin{proposition}[Transfer to an arbitrary odd multiplier]\label{prop:transfer}
Let $m\ge3$ be odd and $a\ge2$. Suppose $H:\UU_a\to\Z$
satisfies \eqref{eq:anti}--\eqref{eq:lipschitz}, with
$G(u)=H(mu)-mH(u)$, the size condition \eqref{eq:size}, and
\[
 \sum_{u\in\UU_a}\wt(G(u))=(m-1)\sum_{u\in\UU_a}\wt(H(u)).
\]
Form the signed digit trees with value functions
\[
 X=1+4H,\qquad Y=-1+4G,
\]
and let $A,B$ be their associated dyadic sign functions. Then
\begin{equation}\label{eq:fexplicit}
 f(n)=\begin{cases}A(n),&m\nmid n,\\B(n/m),&m\mid n\end{cases}
\end{equation}
is a $2^am$-periodic sign function of mean zero and $\LL(f)=0$.
\end{proposition}
\begin{proof}
Antisymmetry and \eqref{eq:mean} give $\sum A=\sum B=0$ over
their dyadic period. Multiplication by odd $m$ permutes all dyadic
layers, so $A\circ m:n\mapsto A(mn)$ also has zero mean and the
same layer sums as $A$.
Put $C=mA+B-A\circ m$. Its value function is
\[
 mX(u)+Y(u)-X(mu)=m-2,
\]
which is constant. The pairing lemma gives
\[
 E(A)=-2^{2-a}\sum_u\wt(H(u)),\qquad
 E(B)= 2^{2-a}\sum_u\wt(G(u)).
\]
Thus $E(C)=(m-1)E(A)+E(B)=0$. Lemma~\ref{lem:cancellation}
implies $\LL(C)=0$.

The function in \eqref{eq:fexplicit} is a sign at every integer,
including multiples of the full period. In the period $2^am$, the
sum of $A(n)$ is $m\sum A=0$, the sum removed at multiples of $m$
is $\sum A(mn)=0$, and the replacement sum is $\sum B=0$.
Hence its mean is zero. Finally,
\[
 f(n)=A(n)+\mathbf1_{m\mid n}\{B(n/m)-A(n)\}
\]
and convergence of all three mean-zero dyadic series gives
\[
 \LL(f)=\LL(A)+\frac{\LL(B)-\LL(A\circ m)}m
         =\frac{\LL(C)}m=0.
\]
This argument uses no primality assumption on $m$.
\end{proof}

\begin{proof}[Proof of Theorem~\ref{thm:main}]
Use the word of Lemma~\ref{lem:word}, place it on the paired cycles
of Lemma~\ref{lem:cycles}, and lift the resulting functions to depth
\eqref{eq:adepth}. The size and valuation requirements of
Proposition~\ref{prop:transfer} all hold. That proposition constructs
$f$ and proves exact vanishing. The depth bound follows from
$J\le2\ell+3$. Finally $2^\ell<2m$, so
\[
 2^am\le2^{5\ell+10}m<2^{15}m^6=32768m^6.
\]
\end{proof}

\section{Exclusions and consequences}\label{sec:consequences}
\begin{proposition}[No pure power of two]\label{prop:pure}
For every $a\ge0$, $2^a\notin\Dset$.
\end{proposition}
\begin{proof}
The case $a=0$ is excluded by the mean condition. Suppose $a\ge1$
and that a periodic sign function $f$ has period $2^a$ and vanishing
harmonic sum. The necessary unit equations in Okada's criterion
\cite[Theorem 1.1]{MP2022}, specialized to the monoid generated by
$2$, give, for every odd integer $u$,
\[
 0=\sum_{j=0}^{a-1}2^{-j}f(2^ju)+2^{1-a}f(0).
\]
The second term is the exact geometric tail for $j\ge a$.
Write $c=f(0)$. Multiplication by $2^{a-1}$ shows
\[
 2^{a-1}f(u)+\sum_{j=1}^{a-1}2^{a-1-j}f(2^ju)+c=0.
\]
The sum of the absolute weights after the first term is exactly
$2^{a-1}$. Equality can hold only when all those signs are opposite
to $f(u)$. Thus $f(u)=-c$ on odd residues, while $f(2^ju)=c$
on every nonzero even layer, as well as at zero. Consequently
$f(n)=c(-1)^n$ and $\LL(f)=-c\log2\ne0$, a contradiction.
\end{proof}

\begin{proof}[Proof of Corollary~\ref{cor:saturation}]
If $n$ is a power of two, all its dyadic multiples are excluded by
Proposition~\ref{prop:pure}. Otherwise write $n=2^bm$ with $m\ge3$
odd, and apply Theorem~\ref{thm:main} to $m$. If its constructed
depth is below $b$, enlarge the period by a power of two. This gives
a dyadic multiple of $n$ in $\Dset$. Closure under multiples also
shows that each set of admissible depths is an upper interval of integers.
\end{proof}

\begin{corollary}[A minimal period for every odd prime]
For each odd prime $p$, $2^{a_0(p)}p$ is a minimal element of
$\Dset$ in the divisibility order. These elements are distinct for
distinct primes, so there are infinitely many such minimal elements.
\end{corollary}
\begin{proof}
A proper divisor either is a power of two, excluded above, or equals
$2^bp$ with $b<a_0(p)$, excluded by definition. Distinct odd prime
factors make the resulting integers distinct.
\end{proof}

This does not assert that the corresponding function has been given
at its least possible depth. It also does not claim that infinitely
many minimal elements were previously unknown: [093] already obtains
an infinite family by another route. The strengthening here is that
every odd prime supports one.

\section{Explicit certificates and independent checks}\label{sec:checks}
The construction is algorithmic, even when the full sign period is
too long to print. A compact certificate consists of the cyclic word,
the depth $J$, representatives of the paired cycles, and the chosen
digit depth $a$. Formula \eqref{eq:fexplicit} then evaluates any term
without listing the full period.

For $m=127$, the parameters are
\[
 \ell=7,\quad s=8,\quad J=17,\quad t=1024,\quad k=256.
\]
The baseline deficit is $D=126976$ and $C=252$. After the $k$ baseline
excursions, the remaining budget is
\[
 B=D-kC=62464=2^{10}+2^{12}+2^{13}+2^{14}+2^{15}.
\]
Use $h_{2j+1}=127+2^M$ for these five exponents, use $h_{2j+1}=-1$
for the other $251$ excursions, and set all other entries to $1$.
On paired cycles modulo $2^{17}$, put $H=\pm2^{17}h$ as in the proof.
The exact maxima are
\[
 \max|H|=4311613440,\qquad \max|G|=547574775808.
\]
Depth $a=41$ satisfies \eqref{eq:size}, proving
\[
 127\cdot2^{41}\in\Dset.
\]
This is a constructive upper bound, not an assertion that $a_0(127)=41$.
The deliberately conservative depth selection in the verifier may be
reduced further without changing the proof.

\begin{center}
\begin{tabular}{rrrrr}
\toprule
$m$ & $J$ & cycle length $t$ & excursions & verified depth $a$\\
\midrule
3&7&32&0&11\\
5&8&64&0&13\\
7&9&64&16&20\\
9&10&128&0&16\\
11&9&128&32&24\\
13&9&128&1&25\\
15&11&128&32&26\\
17&12&256&0&19\\
31&13&256&64&31\\
127&17&1024&256&41\\
\bottomrule
\end{tabular}
\end{center}
None of the displayed depths is claimed optimal. In particular, small
depths already known for $3,5,17$ are much better than these generic outputs.

The integer verifier in Appendix~\ref{app:code} checks every paired
cycle, antisymmetry, the exact valuation budget, signed-digit compatibility,
the constant master identity, both zero period sums, and the vanishing
logarithmic layer statistic. All odd multipliers $3\le m\le63$ and
the additional multiplier $127$ passed. A separate implementation
checked the cyclic budget, order, negation separation, and uniform size
bound for all $4,999$ odd multipliers from $3$ through $9,999$.
It also checked $4,072$ signed-digit pairing identities with a separate
recursive decoder. These are finite checks of formulas; the proofs
above, rather than extrapolation from the checks, establish the infinite theorem.

\section{What the research request has and has not achieved}
The main theorem closes the exact existence conjecture in [093], without
changing its prime quantifier, and strengthens it to all odd multipliers.
It also answers the newly formulated finite cyclic-word problem and
transfers its solution to harmonic cancellation. It does not settle the
external problem of classifying every individual $N\in\Dset$.

The companion results below have distinct statuses. The graph
appendix proves a general finite-core bound, but leaves the unrestricted
Randi\'c--radius conjectures open. The Steklov appendix refutes the
question as printed in the cited 2020 list, even after adding unit
perimeter, but uses classical spectral tools. That exact mathematical
answer does not establish that the observation is historically new.
The randomized-time quantum-walk appendix resolves every higher odd-prime-power
modulus left outside the cited classification; it does not
classify all paths, and it is distinct from standard Ces\`aro averaging.

\begin{center}\small
\begin{tabular}{p{.43\linewidth}p{.48\linewidth}}
\toprule
Requested condition & Supported conclusion in this manuscript\\
\midrule
An unchanged, previously stated target & Yes for [093], Conjecture 9.1; the external period-classification question is retained and remains open here.\\
An unconditional complete proof & Yes for the dyadic-lift theorem and the explicitly stated companion theorems.\\
A complete solution of a longstanding external named problem & The printed Steklov statement is refuted; its continuing open status and priority are not certified. The external number-theory classification is not solved.\\
A historically unprecedented proof principle & Not established. The new cyclic filling step is distinguished from the preceding corpus construction, but its ingredients are classical.\\
No identical prior solution & Targeted source checks found no matching general dyadic-lift theorem; this is not a proof of absence.\\
A result more valuable than every earlier corpus result & No justified ranking is available.\\
\bottomrule
\end{tabular}
\end{center}

Primary sources were inspected on 2 November 2026. Searches included the
Okada-space period problem, dyadic periodic sign cancellation, and the
Steklov common-point question. Mathematical correctness, bibliographic
priority, and comparative significance are separate claims. The strict
conjunction of all requested success conditions has therefore not been
established, and this document should not be represented as such a claim.

\appendix

\section{A finite source budget for supportive graph cores}
\label{sec:graph-source-budget}

This section closes a specific auxiliary gap in corpus manuscript
\cite[Remark~\texttt{rem:D6}]{Corpus087}: for each fixed maximum degree,
the stated supportive-core conditions force an explicit finite bound on
core order. It does not establish the unrestricted Randi\'c--radius
conjectures. The proof combines an exact deletion identity with a forest
whose degrees increase away from its roots.

All graphs in this section are finite and simple. Write
\[
 M(G)=\sum_{uv\in E(G)}\frac{1}{\max(d_G(u),d_G(v))}.
\]
An edge $e$ is \emph{supportive} if $M(G-e)>M(G)$. A
\emph{decorated core} is a graph $G$ obtained from a graph $K$ by attaching
at most one pendant leaf to each vertex of $K$. In this section ``core''
only names the distinguished subgraph: no connectivity hypothesis on
$K$ is needed for the principal theorem. All degrees denoted by $d_v$
are degrees in $G$.

\begin{lemma}[Exact deletion formulas]
\label{lem:graph-exact-deletion}
Let $uv$ be an edge whose endpoints have degrees at least two, and set
\[
 \beta_u=|\{w\in N_G(u)\setminus\{v\}:d_w<d_u\}|,
 \qquad
 \beta_v=|\{w\in N_G(v)\setminus\{u\}:d_w<d_v\}|.
\]
Then
\begin{equation}\label{eq:graph-delete-edge}
 M(G-uv)-M(G)=
 \frac{\beta_u}{d_u(d_u-1)}+
 \frac{\beta_v}{d_v(d_v-1)}-
 \frac1{\max(d_u,d_v)}.
\end{equation}
If $\ell$ is a leaf with support $p$ of degree $d\ge2$, and
$\lambda$ counts the neighbors of $p$ other than $\ell$ whose degrees
are less than $d$, then
\begin{equation}\label{eq:graph-delete-leaf}
 M(G)-M(G-\ell)=\frac1d-\frac{\lambda}{d(d-1)}\ge0.
\end{equation}
\end{lemma}
\begin{proof}
Deleting $uv$ removes its weight. The weight of another incident edge
$uw$ changes only when $d_w<d_u$, in which case its increase is
$1/(d_u-1)-1/d_u=1/[d_u(d_u-1)]$. The corresponding statement at $v$
gives~\eqref{eq:graph-delete-edge}. For leaf deletion, the removed
weight is $1/d$ and the other affected edges have exactly the same
increments. Since $\lambda\le d-1$, the resulting drop is nonnegative.
These are the elementary degree calculations used in
\cite{Corpus087,Corpus094}.
\end{proof}

\begin{theorem}[Source budget and finite core bound]
\label{thm:graph-source-budget}
Let $D\ge3$ be an integer and let $G$ be a decorated core with
\[
 3\le d_v\le D\quad(v\in V(K)),
 \qquad M(G-e)>M(G)\quad(e\in E(K)).
\]
If $M(G)-M(K)<1$, then
\begin{equation}\label{eq:graph-factorial-bound}
 |V(K)|\le B(D):=3\sum_{j=1}^{D-2}j!-1.
\end{equation}
\end{theorem}
\begin{proof}
Call $s\in V(K)$ a \emph{source} if it has no core neighbor of strictly
smaller full degree, and let $S$ be the set of sources.

Every source has a pendant leaf. Otherwise a source $u$ would have no
lower-degree neighbor in $G$. For any core edge $uv$ incident to $u$,
we would have $d_v\ge d_u$, $\beta_u=0$, and $\beta_v\le d_v-1$.
Formula~\eqref{eq:graph-delete-edge} would give
\[
 M(G-uv)-M(G)
 \le\frac{d_v-1}{d_v(d_v-1)}-\frac1{d_v}=0,
\]
contrary to supportiveness. Such an edge exists because $d_u\ge3$
and there is at most one pendant neighbor.

The set $S$ is independent. If adjacent sources have unequal degrees,
the higher-degree endpoint has a lower-degree core neighbor, which
contradicts its definition. If adjacent sources $u,v$ both have degree
$d$, their only lower-degree neighbors are their respective leaves.
Thus $\beta_u=\beta_v=1$, and
\[
 M(G-uv)-M(G)=\frac2{d(d-1)}-\frac1d
             =\frac{3-d}{d(d-1)}\le0.
\]
The equality case $d=3$ also fails the strict supportive condition.

Let $L_S$ be the leaves attached to sources. Deleting all of $L_S$
changes no surviving edge weight. Indeed, sources are independent,
and every core neighbor of a source $s$ has unchanged degree at least
the original $d_s$. Consequently
\begin{equation}\label{eq:graph-source-identity}
 M(G)-M(G-L_S)=\sum_{s\in S}\frac1{d_s}.
\end{equation}
Delete the remaining pendant leaves one at a time. Their supports still
have degree at least two, because their original core degrees were at
least two. Each deletion has nonnegative drop by
\eqref{eq:graph-delete-leaf}. Therefore
\begin{equation}\label{eq:graph-source-budget}
 \sum_{s\in S}\frac1{d_s}
 \le M(G)-M(K)<1.
\end{equation}

For every non-source core vertex choose one lower-degree core neighbor
as its parent. Parent degrees strictly decrease, so the chosen edges
form a rooted forest with root set $S$. Every root of degree $d$ has at
most $d-1$ children, because one incident edge leads to its pendant
leaf. Every non-root of degree $d$ has at most $d-1$ children, because
one incident edge leads to its parent. Degrees strictly increase along
each root-to-leaf path.

Define integers
\begin{equation}\label{eq:graph-forest-recursion}
 F_D(D)=1,\qquad
 F_D(d)=1+(d-1)F_D(d+1)\quad(3\le d<D).
\end{equation}
They decrease as $d$ increases. Induction on $D-d$ therefore bounds a
subtree rooted at a degree-$d$ vertex by $F_D(d)$, even if some children
skip degree levels. Hence
\[
 |V(K)|\le\sum_{s\in S}F_D(d_s).
\]
For $3\le d<D$ the recurrence gives
\[
 dF_D(d)-(d+1)F_D(d+1)
 =d+(d^2-2d-1)F_D(d+1)>0.
\]
Thus $dF_D(d)\le3F_D(3)$ for every admissible degree. By
\eqref{eq:graph-source-budget},
\[
 |V(K)|\le3F_D(3)\sum_{s\in S}\frac1{d_s}<3F_D(3).
\]
Expanding~\eqref{eq:graph-forest-recursion} gives
$F_D(3)=\sum_{j=1}^{D-2}j!$. Integrality proves
\eqref{eq:graph-factorial-bound}.
\end{proof}

\begin{proposition}[Exact integer refinement]
\label{prop:graph-knapsack}
Under the hypotheses of Theorem~\ref{thm:graph-source-budget}, the bound
$B(D)$ can be replaced by
\begin{equation}\label{eq:graph-knapsack}
 B_*(D)=\max\left\{
 \sum_{d=3}^D n_dF_D(d):
 n_d\in\mathbb Z_{\ge0},\quad
 \sum_{d=3}^D\frac{n_d}{d}<1
 \right\}.
\end{equation}
The following values are exact:
\begin{center}
\begin{tabular}{r|rrrrr}
$D$&4&5&6&7&8\\\hline
$B(D)$&8&26&98&458&2618\\
$B_*(D)$&7&22&82&382&2182
\end{tabular}
\end{center}
These are bounds from the source budget, not claims that corresponding
graphs exist.
\end{proposition}
\begin{proof}
Take $n_d$ to be the number of degree-$d$ sources in the proof above.
The feasible set in~\eqref{eq:graph-knapsack} is finite, because
$n_d\le d-1$. The table is reproduced by the following exact recurrence.
Put $L=\operatorname{lcm}(3,\ldots,D)$ and initialize $A(0)=0$ and
$A(w)=-\infty$ for $1\le w<L$. Process $w=0,1,\ldots,L-1$ in order;
for each $d$ with $w+L/d<L$, update
\[
 A(w+L/d)\ \leftarrow\
 \max\{A(w+L/d),\ A(w)+F_D(d)\}.
\]
The maximum of the final $A(w)$ is $B_*(D)$. Positive integer weights
ensure that this recurrence visits every admissible multiset of source
degrees. Only integer arithmetic is involved.
\end{proof}

\subsection{Application conditional on the corpus reduction}

Let $\rho(G)$ be the absolute radius: minimize the maximum distance to
a vertex over all points of the unit-edge metric realization of $G$.
Equivalently, $2\rho(G)$ is the minimum diameter of a spanning tree.
The stronger inequality $M(G)\ge\rho(G)$ would imply
$R(G)\ge\operatorname{rad}(G)-\tfrac12$, and hence Fajtlowicz's
$R(G)\ge\operatorname{rad}(G)-1$. The chemical-graph results of
\cite{CPS} provide established context for this approach.

\begin{corollary}[Finite criterion using the earlier structure theorem]
\label{cor:graph-finite-criterion}
Assume the minimal-counterexample reduction proved in
\cite[Proposition~\texttt{prop:structure}]{Corpus087}. For each fixed
$D\ge3$, the following are equivalent:
\begin{enumerate}
\item $M(G)\ge\rho(G)$ for every connected graph with maximum degree
at most $D$;
\item the same inequality holds for every such graph with at most
$2B_*(D)$ vertices.
\end{enumerate}
One may replace $2B_*(D)$ by the closed bound $2B(D)$.
\end{corollary}
\begin{proof}
Only the second implication needs proof. If the universal assertion
fails, the cited reduction supplies a minimal counterexample $G=K+L$
with at most one leaf per core vertex, all full core degrees between
three and $D$, every core edge supportive, and $M(G)-M(K)<1$.
Proposition~\ref{prop:graph-knapsack} bounds $|K|$ by $B_*(D)$.
Since $|L|\le|K|$, this contradicts the second assertion.
\end{proof}

The external dependency in this corollary is explicit; the present
section does not reproduce the earlier cut-vertex padding and
absolute-radius arguments. For $D=6$ the criterion requires checking
orders at most $164$, or $196$ with the simpler factorial bound.
Neither check has been performed here. The finite reduction therefore
settles the auxiliary fixed-degree size question, while leaving both
the degree-six inequality and the unrestricted named conjectures
unresolved. It also supplies no claim of worldwide priority for the
source-budget method.

\subsection{Independent exact-arithmetic audit}

A separate checker enumerated every labelled core on $3$ through $6$
vertices with minimum core degree at least two, and every subset of
vertices receiving a pendant leaf. It recomputed $M(G-e)$ directly
from each edge-deleted graph using exact rational arithmetic, rather
than using~\eqref{eq:graph-delete-edge} to decide supportiveness.
Among $780\,616$ decorated graphs, $736$ had all core edges supportive
and all full core degrees at least three. For every accepted graph,
the checker verified source independence, the exact identity
\eqref{eq:graph-source-identity}, its lower bound on total leaf drop,
and the forest estimate. It also checked
\eqref{eq:graph-factorial-bound} whenever the total drop was less than
one. Every check passed. These computations check small cases and the
implementation; the all-orders result rests on the analytic proof.
\section{A separate spectral audit: common small values of Steklov eigenfunctions}
\label{sec:steklov-audit}

This section records an independent, elementary correction to a question
in the April 2020 eigenvalue miniconference list. It is separate from the
arithmetic construction in the main text. Open Problem~10, attributed to
Jean Lagac\'e in \cite{Lagace2020}, asks whether a surface or planar domain
has a point at which every first positive Steklov eigenfunction, normalized
in boundary $L^2$, has absolute value below $1/\sqrt2$, or below some other
universal constant. The printed formulation imposes neither simple
connectivity nor a boundary-length normalization. We address that printed
formulation; no claim is made about an unstated intended restriction.

The proposed threshold fails even after imposing perimeter one. Without
scale normalization, no universal positive constant works. A natural
replacement involving eigenspace dimension is valid, and its constant is
optimal when the first eigenspace has dimension two. All ingredients are
classical: separation of variables, finite-dimensional evaluation norms,
and the subharmonicity of a harmonic vector norm. In particular, this
section does not claim a historically new proof principle. The specific
priority of the correction below has not been certified.

\subsection{Evaluation norms and a universal strict bound}

Let $\Omega\subset\R^d$ be a bounded connected Lipschitz domain, with
boundary measure $L=|\partial\Omega|$. Let $H$ be an $m$-dimensional real
space of harmonic functions in $H^1(\Omega)$ whose boundary traces are
linearly independent. Give $H$ the boundary $L^2$ norm, and choose an
orthonormal basis $u_1,\ldots,u_m$. At an interior point define
\[
 \kappa_H(x)=\sup\{|u(x)|:u\in H,\ \|u\|_{L^2(\partial\Omega)}=1\}.
\]
Cauchy--Schwarz, with equality for a coefficient vector proportional to
$(u_1(x),\ldots,u_m(x))$ when this vector is nonzero, gives
\begin{equation}\label{eq:steklov-evaluation}
 \kappa_H(x)^2=\sum_{j=1}^m u_j(x)^2.
\end{equation}
Thus this expression is independent of the orthonormal basis.

\begin{proposition}[A common small-value point]\label{prop:steklov-universal}
If $H$ contains a nonconstant function, there is an interior point $x$ such
that
\[
 \kappa_H(x)<\sqrt{m/L}.
\]
In particular, this holds for every positive Steklov eigenspace, and for
any finite direct sum of positive Steklov eigenspaces, with $m$ its total
dimension.
\end{proposition}

\begin{proof}
Put $U=(u_1,\ldots,u_m)$ and $w=|U|$. The Lipschitz chain rule gives
$w\in H^1(\Omega)$. The function $w$ is weakly subharmonic: locally the
smooth approximants $(\epsilon^2+|U|^2)^{1/2}$ have nonnegative Laplacian
by Cauchy--Schwarz, and their gradients converge locally in $L^2$ to
$\nabla w$. Consequently
\[
 \int_\Omega\nabla w\cdot\nabla\varphi\leq0
 \qquad(\varphi\in H^1_0(\Omega),\ \varphi\geq0),
\]
first for smooth compactly supported tests and then by positive
approximation in $H^1_0$.

Suppose $w\geq c:=\sqrt{m/L}$ throughout the interior. The Sobolev trace
preserves order and commutes with the norm map. Orthonormality therefore
implies
\[
 \operatorname{Tr}w\geq c,\qquad
 \int_{\partial\Omega}(\operatorname{Tr}w)^2=m=c^2L.
\]
Hence $\operatorname{Tr}w=c$ almost everywhere, so $w-c$ is a nonnegative
member of $H^1_0(\Omega)$. Testing subharmonicity with $w-c$ yields
$\int_\Omega|\nabla w|^2\leq0$. Thus $w$ is constant. Interior smoothness
of the harmonic functions now gives
\[
 0=\Delta|U|^2=2\sum_{j=1}^m|\nabla u_j|^2,
\]
forcing every member of $H$ to be constant, a contradiction.
For a positive Steklov eigenvalue a nonzero constant cannot satisfy
$\partial_\nu u=\sigma u$, proving the stated application.
\end{proof}

\subsection{The first eigenspace of a round annulus}

For $0<a<1$ let $A_a=\{x\in\R^2:a<|x|<1\}$. We use the usual convention
$\Delta u=0$ in $A_a$ and $\partial_\nu u=\sigma u$ on both boundary
circles. The separation of variables recalled below is classical; see
\cite[Example~4.2.5]{GP2017}.

\begin{lemma}\label{lem:steklov-annulus}
The first positive Steklov eigenvalue of $A_a$ has multiplicity exactly
two. Its eigenspace is
\[
 \operatorname{span}\{F_a(r)\cos\theta,F_a(r)\sin\theta\},
 \qquad F_a(r)=r+c_a/r,
\]
where $c_a$ is the positive root of
\begin{equation}\label{eq:steklov-c-root}
 c_a^2+\frac{(1-a)(1+a^2)}{1+a}c_a-a^2=0,
 \qquad a^2<c_a<a.
\end{equation}
The eigenvalue is $(1-c_a)/(1+c_a)$.
\end{lemma}

\begin{proof}
The angular frequency $n\geq1$ has harmonic radial profiles
$Ar^n+Br^{-n}$. Its lowest eigenvalue is the minimum of
\[
 Q_n(f)=
 \frac{\displaystyle\int_a^1\left(r|f'|^2+n^2|f|^2/r\right)\,dr}
 {\displaystyle |f(1)|^2+a|f(a)|^2}.
\]
For every nonzero admissible profile, $Q_n(f)>Q_1(f)$ when $n\geq2$.
Applying this to a minimizing profile in the $n$th sector shows that its
lowest eigenvalue strictly exceeds that in the first sector. The radial
sector $n=0$ consists of the zero eigenvalue and the eigenvalue
\[
 \sigma_{\mathrm{rad}}=\frac{1+a}{a\log(1/a)}>1;
\]
the last inequality follows from $\log(1/a)<a^{-1}-1$.

For frequency one write $F=r+c/r$. The two boundary equations imply
$c=(1-\sigma)/(1+\sigma)$ and
\[
 a(1-a^2)\sigma^2-(1+a)(1+a^2)\sigma+(1-a^2)=0.
\]
This quadratic is positive at zero and negative at one. Its leading
coefficient is positive, so its smaller root lies in $(0,1)$ and the
other exceeds one. Hence the smaller root is the first positive
eigenvalue of the whole annulus, occurring only in the sine and cosine
modes. Eliminating $\sigma$ gives \eqref{eq:steklov-c-root}. The polynomial
there is negative at $c=a^2$ and positive at $c=a$; its unique positive
root therefore satisfies the displayed bounds.
\end{proof}

The two basis functions in Lemma~\ref{lem:steklov-annulus} are orthogonal
on the boundary and have common squared norm $\pi D_a$, where
\[
 D_a=(1+c_a)^2+a(a+c_a/a)^2.
\]
Because $a<\sqrt{c_a}<1$, the minimum of $F_a$ occurs at
$r=\sqrt{c_a}$ and equals $2\sqrt{c_a}$. Scaling a planar domain by a
factor $s$ multiplies boundary measure by $s$ and the squared evaluation
norm of boundary-unit eigenfunctions by $s^{-1}$. Thus the annulus
\[
 \widehat A_a=\frac1{2\pi(1+a)}A_a,
 \qquad |\partial\widehat A_a|=1,
\]
has minimum squared first-eigenspace evaluation norm
\begin{equation}\label{eq:steklov-annular-M}
 M(a)=\min_{x\in\overline{\widehat A_a}}\kappa_{E_1}(x)^2
     =\frac{8(1+a)c_a}{(1+c_a)^2+a(a+c_a/a)^2}.
\end{equation}
The minimum occurs in the interior, so using the closure here does not
change its value.

\subsection{An exact counterexample and the optimal replacement}

\begin{theorem}[An exact perimeter-one counterexample]
\label{thm:steklov-counterexample}
For the smooth connected domain
\[
 \Omega=\frac1{3\pi}\{x\in\R^2:\tfrac12<|x|<1\},
 \qquad |\partial\Omega|=1,
\]
the first positive Steklov eigenspace satisfies
\[
 \min_{x\in\overline\Omega}\kappa_{E_1}(x)^2
 =-8+\frac{40\sqrt{17}}{17}
 =\frac{16(\sqrt{17}-1)}{17+3\sqrt{17}}>\frac12.
\]
Consequently no point of $\Omega$ has the common small-value property
with threshold $1/\sqrt2$ posed in the printed question.
\end{theorem}

\begin{proof}
At $a=1/2$, the first eigenvalue on the unscaled annulus is
$(5-\sqrt{17})/2$, and
\[
 c_{1/2}=\frac{\sqrt{17}-1}{8},\qquad
 D_{1/2}=\frac{51+9\sqrt{17}}{32}.
\]
Substitution into \eqref{eq:steklov-annular-M} gives the formula. For a
simple exact strict comparison, $4<\sqrt{17}<5$ makes its last expression
larger than $16\cdot3/32=3/2$. At each point, the supremum defining
$\kappa_{E_1}$ is attained by a unit vector of coefficients. Thus the
failure is of the required universal quantifier over all normalized
eigenfunctions, rather than merely a comparison of one chosen basis.
\end{proof}

\begin{corollary}[Scaling and sharp multiplicity-two correction]
\label{cor:steklov-sharp}
Without fixing a geometric scale, no positive universal constant can
replace $1/\sqrt2$ in the printed question. Among smooth connected planar
domains whose first positive Steklov eigenspace has dimension two, the
least universal constant $C$ such that some interior point satisfies
\[
 |u(x)|<\frac{C}{\sqrt{|\partial\Omega|}}
 \quad\text{for every }u\in E_1(\Omega),\quad
 \|u\|_{L^2(\partial\Omega)}=1,
\]
is exactly $C=\sqrt2$.
\end{corollary}

\begin{proof}
The strictly positive minimum in Theorem~\ref{thm:steklov-counterexample}
becomes arbitrarily large under shrinking dilations, proving the first
assertion. Proposition~\ref{prop:steklov-universal} proves the second
assertion with $C=\sqrt2$. Conversely, \eqref{eq:steklov-c-root} gives
$c_a\to1$ as $a\uparrow1$, and hence $M(a)\to2$. Every smaller constant
therefore fails on a sufficiently thin round annulus of perimeter one.
For consistency with the strict universal inequality, note also that
\[
 D_a-4(1+a)c_a=(1-c_a)^2+a(a-c_a/a)^2>0,
\]
so $M(a)<2$ for every individual annulus. The sharp constant is approached
but the limiting thin annulus is not itself a domain in the class.
\end{proof}

The last corollary is an optimal statement for the explicitly specified
multiplicity-two class. It does not assert that the first positive
Steklov eigenvalue always has multiplicity two, nor does it settle
additional topological restrictions absent from the printed question.
The original problem list establishes the provenance of the question;
the classical annulus formula establishes the provenance of the main
computational ingredient. Neither fact certifies historical priority of
this particular correction.
\section{Exact randomized uniform mixing on ternary-power paths}
\label{sec:ternary-paths}

This section concerns a random observation time for the ordinary adjacency
quantum walk. It does not concern the standard Ces\`aro average mixing
matrix. In particular, its conclusion does not answer the rank-two
Ces\`aro question discussed in corpus item [077].

Coutinho and Guo explicitly asked whether the randomized uniform mixing
known for $P_3$ and $P_4$ extends to larger paths
\cite[Section~4.2]{CoutinhoGuo2024}. Kang subsequently proved positive
results when the number of vertices plus one is a prime, twice a prime,
or a power of two, and several infinite families of negative results
\cite{Kang2026}. The latter paper explicitly leaves the moduli outside
its stated positive and negative families unresolved. Thus the theorem
below resolves an infinite surviving subfamily; it neither gives an
all-path classification nor asserts that every path has the property.

\begin{theorem}[Ternary-power paths]
\label{thm:ternary-uam}
Let $k\ge2$ and $n=3^k-1$. If $A$ is the adjacency matrix of the
unweighted path $P_n$, there is a bounded, compactly supported probability
density $\rho$ on $[0,\infty)$ such that
\[
 \int_0^\infty |(e^{itA})_{uv}|^2\rho(t)\,dt=\frac1n
 \qquad(1\le u,v\le n).
\]
The same density works simultaneously for every initial vertex.
\end{theorem}

\subsection{The covariance that yields uniform mixing}
Put $q=n+1$. The eigenvalues and normalized eigenvectors are
\[
 \lambda_j=2\cos\frac{\pi j}{q},\qquad
 v_j(u)=\sqrt{\frac2q}\sin\frac{\pi ju}{q}
 \quad(1\le j,u\le q-1).
\]
Write $E_j=v_jv_j^{\mathsf T}$, let $R$ reverse the vertex order, and
let $J$ be the all-ones matrix. The sine product identity gives
\begin{equation}\label{eq:ternary-cesaro}
 D_0:=\sum_j E_j\circ E_j
   =\frac1q\left(J+\frac{I+R}{2}\right).
\end{equation}
Indeed, expand the product of two squared sines into cosines and use
\(\sum_{j=1}^{q-1}\cos(2\pi aj/q)=q\mathbf1_{q\mid a}-1\).
This proves the displayed identity entry by entry.
Furthermore $Rv_j=(-1)^{j+1}v_j$, so
\begin{equation}\label{eq:ternary-parity}
 \sum_{j\equiv\ell\ ({\rm mod}\ 2)}E_j\circ E_\ell
 =\frac12\left[
   \left(\sum_jE_j\right)\circ\left(\sum_jE_j\right)
   +\left(\sum_j(-1)^jE_j\right)\circ
    \left(\sum_j(-1)^jE_j\right)\right]
 =\frac{I+R}{2}.
\end{equation}

Suppose phases $z_1,\ldots,z_n$ have covariance
\begin{equation}\label{eq:ternary-covariance}
 \mathbb E[z_j\overline z_\ell]=
 \begin{cases}
  1,&j=\ell,\\
  -c,&j\ne\ell\text{ and }j\equiv\ell\pmod2,\\
  0,&j\not\equiv\ell\pmod2.
 \end{cases}
\end{equation}
The expected mixing matrix is then
\[
 (1+c)D_0-c\frac{I+R}{2}.
\]
At $c=c_0:=1/n$ this is exactly $J/n$. Since $q$ is odd,
each pair $j,q-j$ has one odd and one even member. It therefore suffices
to construct odd-index phases $W_1,\ldots,W_m$, $m=n/2$, satisfying
\begin{equation}\label{eq:ternary-odd-covariance}
 \mathbb E[W_i\overline W_j]=-c\quad(i\ne j),
 \qquad \mathbb E[W_iW_j]=0\quad\text{for all }i,j,
\end{equation}
and set the even-index phases to the conjugates of their reflected
odd-index partners. These are precisely the conjugacy constraints of
the actual time flow, because $\lambda_{q-j}=-\lambda_j$.

\subsection{All frequency relations, including the integer lattice}
Write $q=3s$, where $s=3^{k-1}$, and put
\[
 r=\frac{s-1}{2},\qquad m=3r+1,\qquad n=6r+2.
\]
The odd indices partition into the singleton $s$ and the disjoint triples
\[
 \{j,2s-j,2s+j\},\qquad j=1,3,\ldots,s-2.
\]
Here $\lambda_s=1$, and the cosine addition formula gives
\begin{equation}\label{eq:ternary-triples}
 \lambda_j+\lambda_{2s-j}+\lambda_{2s+j}=0.
\end{equation}
There are no further independent rational relations. To prove this,
the real cyclotomic field
\(K=\mathbb Q(2\cos(\pi/q))\) has degree
\(\varphi(2q)/2=s\). The odd-mode frequencies span $K$ over
$\mathbb Q$: powers of $\lambda_1$ are combinations of $1$ and
$2\cos(a\pi/q)$, the latter reduce by periodicity and reflection to
the displayed modes, every even mode equals the negative of an odd
mode, and $1=\lambda_s$ is already present. Thus their span has
dimension $s=m-r$. The $r$ disjoint relations
\eqref{eq:ternary-triples} are independent and exhaust the rational
relation space. They also generate the integer relation lattice:
in any rational combination of the disjoint relation vectors, the
three coordinates on each triple equal its combination coefficient;
integral coordinates force that coefficient to be integral.

It follows from the character criterion for a linear torus flow that
\(t\mapsto(e^{it\lambda_j})_{j\ {
m odd}}\) is dense and
equidistributed in exactly the torus
\begin{equation}\label{eq:ternary-torus}
 \mathcal T=\left\{(W_0,(W_{a1},W_{a2},W_{a3})_{a=1}^r):
 |W_i|=1,\quad W_{a1}W_{a2}W_{a3}=1\right\}.
\end{equation}
Its dimension is $2r+1=s$. The singleton $W_0$ is unrestricted.
This step accounts for the dependent modes; it does not replace them
with independent phases.

\subsection{An exchangeable constrained-phase construction}
First suppose $r>1$. We construct \eqref{eq:ternary-odd-covariance}
at a value $c>c_0$. Set explicitly
\[
 c=\frac1{4(r+1)},\qquad
 B=\sqrt{3-6c},\qquad
 \alpha=\frac{9c}{B^2}=\frac3{4r+2},\qquad \beta=\frac B3.
\]
Then $c>c_0$, $0<\beta<1$, and
\(\alpha<1/(r+1)\). Let $L=r+1$ and
$\omega=e^{2\pi i/3}$. We need exchangeable cube-root labels
$Y_0,\ldots,Y_r$ with
\(\mathbb E[Y_a\overline Y_b]=-\alpha\) for distinct positions.

Distribute $L$ labels among $1,\omega,\omega^2$ with counts differing
by at most one, uniformly permute the positions, and multiply all labels
by an independent uniform cube root. If $\epsilon=0$ when $3\mid L$
and $\epsilon=1$ otherwise, the squared modulus of their sum is
$\epsilon$. Hence their common off-diagonal correlation is
\[
 -b_L=\frac{\epsilon-L}{L(L-1)},\qquad
 b_L\ge\frac1L.
\]
Mixing this distribution with independent uniform cube-root labels
realizes every correlation between $-b_L$ and zero, including $-\alpha$.

Choose $\theta$ with $\cos\theta=(B-1)/2$. Conditional on the labels,
independently and uniformly permute the three entries of
\[
 (W_{a1},W_{a2},W_{a3})
   =Y_a(1,e^{i\theta},e^{-i\theta}),\qquad 1\le a\le r.
\]
Their product is one and their conditional sum is $BY_a$. Set the
singleton $W_0=Y_0$ with probability $\beta$ and otherwise take an
independent uniform cube root. A distinct pair in the same triple has
correlation
\[
 \frac{B^2-3}{6}=-c.
\]
A pair from different triples has correlation
\(B^2(-\alpha)/9=-c\), while a singleton--triple pair has correlation
\((B/3)\beta(-\alpha)=-c\). The law is invariant under simultaneous
multiplication of every phase by $\omega$, which preserves every
triple-product constraint. Therefore
\(\mathbb E[W_iW_j]=\omega^2\mathbb E[W_iW_j]=0\).
This proves \eqref{eq:ternary-odd-covariance} at the chosen $c>c_0$.

For $r=1$, which is $P_8$, use $c=1/6>c_0=1/8$ and $B=\sqrt2$.
Uniformly permute $(1,e^{i\theta},e^{-i\theta})$, where
$\cos\theta=(B-1)/2$, for the single triple. Independently take the
singleton in $\{1,-1\}$ with mean $-3c/B=-1/(2\sqrt2)$.
Apply a common uniform cube-root rotation to all four phases. The
within-triple correlation is $-c$; the singleton--triple correlation
is $(-3c/B)(B/3)=-c$; and the holomorphic second moments vanish
by the same rotation argument. This handles the smallest case with
strict margin as well.

\subsection{Exact transfer to a compactly supported time density}
Let $\nu_c$ denote the constructed measure on $\mathcal T$, and let
$H$ be Haar probability measure on that torus. Every off-diagonal
Hermitian moment and every holomorphic second moment appearing in
\eqref{eq:ternary-odd-covariance} is a nontrivial character: distinct
odd modes have neither equal nor opposite frequencies. Its Haar
average is therefore zero. The mixture
\[
 \mu=\left(1-\frac{c_0}{c}\right)H+\frac{c_0}{c}\nu_c
\]
has the target covariance for $c_0$, and has a strictly positive Haar
component $\eta=1-c_0/c$.

A measure on the orbit closure alone would not yet prove the theorem.
We give the exact finite-time transfer. Symmetrize $\mu$ under inversion,
which preserves the real target moments. Choose the finite list
$F_0=1,F_1,\ldots,F_D$ of distinct cosine characters, modulo sign,
that occur in the mixing entries. Distinct entries of this list are
Haar-orthogonal, with squared norms $1,1/2,\ldots,1/2$.
Convolve $\mu$ with a positive product Fej\'er approximate identity
on $\mathcal T$. Its density is at least $\eta$ and its finitely
many target moments converge to those of $\mu$. If the discrepancy
in moment $F_j$ is $\delta_j$, add $2\sum_{j=1}^D\delta_jF_j$.
Orthogonality corrects every moment exactly, preserves normalization,
and changes the sup norm by at most $2\sum_j|\delta_j|$.
For a sufficiently close approximation we thus obtain a finite
trigonometric-polynomial density $f\ge\eta/2$ with exact target
moment vector $b$ (including $b_0=1$).

Write $F(t)$ and $f(t)$ for their evaluations along the time flow.
Define
\[
 G_T=\frac1T\int_0^T F(t)F(t)^{\mathsf T}\,dt,
 \qquad b_T=\frac1T\int_0^T F(t)f(t)\,dt.
\]
Equidistribution gives
\(G_T\to\operatorname{diag}(1,1/2,\ldots,1/2)\) and
$b_T\to b$. Hence $G_T$ is invertible for all sufficiently large
finite $T$, and
\[
 a_T=G_T^{-1}(b-b_T)\longrightarrow0.
\]
For such a sufficiently large $T$, set
\[
 \rho_T(t)=\frac1T\left(f(t)+a_T^{\mathsf T}F(t)\right)
 \mathbf1_{[0,T]}(t).
\]
The strict lower bound on $f$ and $a_T\to0$ give $\rho_T\ge0$.
The defining linear system gives exactly
\(\int\rho_TF=b\), including \(\int\rho_T=1\).
Thus $\rho_T$ is a bounded compactly supported probability density
on actual observation times, with the required exact moments.
Equations \eqref{eq:ternary-cesaro}--\eqref{eq:ternary-covariance}
now give $J/n$ and finish the proof of
Theorem~\ref{thm:ternary-uam}.

\paragraph{Contribution boundary.}
The theorem covers $P_8,P_{26},P_{80},P_{242},\ldots$.
The cyclotomic relation analysis and constrained-phase construction are
the added ingredients relative to the independent-mode cases in
\cite{Kang2026}. The spectral covariance target, torus methods,
and finite-moment correction have established antecedents and are
not claimed as historically unprecedented principles. A targeted
primary-source review through 2 November 2026 found no earlier
identical ternary-power theorem; this is a bounded priority assessment,
not a certification of worldwide absence. The larger classification
problem remains open outside the established families.
\subsection{Extension to every odd prime power}
\label{sec:odd-prime-power-uam}

The preceding phase construction admits a stronger form. This closes
every odd-prime-power modulus, while retaining the distinction between
randomized observation and the standard Ces\`aro average.

\begin{theorem}[All odd prime powers]
\label{thm:prime-power-uam}
For every odd prime $p$ and every integer $k\ge2$, the path
$P_{p^k-1}$ admits a bounded, compactly supported probability density
on $[0,\infty)$ whose averaged quantum mixing matrix is exactly
$J/(p^k-1)$.
\end{theorem}

\begin{proof}
Put
\[
 q=p^k=ps,\quad h=\frac{p-1}{2},\quad
 r=\frac{s-1}{2},\quad m=\frac{q-1}{2}=pr+h,
 \quad c_0=\frac1{2m}.
\]
We reuse the spectral covariance calculation and exact finite-time
transfer above. It remains to identify the phase torus and construct
the target covariance with a strictly positive Haar component.

\emph{The exact relation lattice.}
For every odd $j$ with $1\le j<s$, consider the exponents
$j+2s\ell$, $0\le\ell<p$, modulo $2q$, and reflect an exponent
larger than $q$ to $2q$ minus that exponent. They give $p$ distinct
odd mode indices in $\{1,\ldots,q-2\}$. Their frequencies have
sum zero by summing a regular $p$-gon. As $j$ ranges over these
$r$ values, the resulting $p$-tuples are disjoint. The remaining
$h$ mode indices are
\[
 s,3s,\ldots,(p-2)s.
\]
One way to verify the partition is to classify odd exponents modulo
$2s$, up to negation. The residues $j$ and $-j$, with odd $j<s$,
give the nonexceptional pairs of classes. The exceptional class is
$s\equiv-s\pmod{2s}$; after reflection and exclusion of the
non-mode index $q$, it gives exactly the displayed $h$ singletons.

The odd-mode frequencies span the real cyclotomic field of degree
\(\varphi(2q)/2=h s\). The proof by powers of
$2\cos(\pi/q)$ is the same as before; the constant is in the span
because
\[
 \sum_{a=1,3,\ldots,p-2}2\cos\frac{\pi a}{p}=1.
\]
Since $m-r=hs$, the disjoint $p$-tuple sums exhaust all rational
relations. Their integer lattice is saturated for the same coordinate
reason as in the ternary case. The actual time-flow closure is
therefore the torus consisting of $r$ blocks of $p$ phases, each
with product one, and $h$ unrestricted singleton phases.

\emph{Exchangeable root labels with a strict margin.}
Assume first $r\ge2$. Let $\omega=e^{2\pi i/p}$.
As $r=pu+h$ for an integer $u\ge0$, form a multiset of $r$ roots
containing $u$ complete sets of the $p$th roots and the $h$ extra
roots
\[
 1,\omega^2,\omega^4,\ldots,\omega^{2h-2}.
\]
Its sum has squared modulus
\[
 d_p=\left|\frac{1-\omega^{2h}}{1-\omega^2}\right|^2
     =\frac1{4\cos^2(\pi/p)}\le1.
\]
Uniformly permute its positions and apply a common independent
uniform $p$th-root rotation. The resulting exchangeable vector
$Y=(Y_1,\ldots,Y_r)$ has zero coordinate means and common
off-diagonal correlation $-b$, where
\[
 b=\frac{r-d_p}{r(r-1)}>\frac1{2r}.
\]
The strict inequality follows from $d_p\le1<(r+1)/2$.
By mixing with independent uniform $p$th-root coordinates, we can
therefore arrange any pair correlation $-\alpha$ with $0\le\alpha\le b$.

Choose $c>c_0$ sufficiently close to $c_0$, and define
\[
 B=\sqrt{p\bigl(1-(p-1)c\bigr)},\qquad
 \alpha=\frac{p^2c}{B^2}=\frac{pc}{1-(p-1)c},\qquad
 \beta=\frac Bp.
\]
At $c=c_0$ we have $\alpha=1/(2r)$. Hence the choice can be made so
that $0<\alpha<b$, $\alpha<1/(2r-1)$, $B>0$, and $0<\beta<1$.
Take $Y$ with pair correlation $-\alpha$ as just constructed.

We also need a second vector $Y'$ with the same internal correlations
and with every cross correlation equal to $-\alpha$. With probability
\[
 t=\frac{\alpha r}{1-(r-1)\alpha}\in(0,1),
\]
let $Y'$ be the negative of a uniform random permutation of $Y$;
otherwise let it be an independent copy of $Y$. Both alternatives
have internal pair correlation $-\alpha$. In the first alternative a
cross correlation is
\[
 -\frac{1-(r-1)\alpha}{r};
\]
in the second it is zero, because the coordinate means vanish.
Thus every cross correlation is $-\alpha$. The coordinates of $Y'$
need only be unit phases: they will be used for the unrestricted
singletons, so the possible minus sign causes no product constraint.
The joint law of $(Y,Y')$ is invariant under common multiplication
by $\omega$.

\emph{Blocks and singletons.}
Choose $\theta$ so that $1+(p-1)\cos\theta=B$. This is possible:
for the chosen $c$ we have $0<B<\sqrt p<p$ and hence
$-1<(B-1)/(p-1)<1$. For each block $a$, independently permute
the $p$ phases
\[
 Y_a\bigl(1,
 \underbrace{e^{i\theta},e^{-i\theta},\ldots,
 e^{i\theta},e^{-i\theta}}_{h\text{ conjugate pairs}}\bigr).
\]
Their product is $Y_a^p=1$ and their conditional sum is $BY_a$.
For singleton $b\in\{1,\ldots,h\}$, retain $Y'_b$ with
probability $\beta$, and otherwise use an independent uniform
$p$th root. Make these choices independently. There are enough
coordinates since $h\le r$.

The correlation of distinct phases within one block is
\[
 \frac{B^2-p}{p(p-1)}=-c.
\]
For phases in different blocks it is
$-(B/p)^2\alpha=-c$. For a singleton and a block phase it is
$-\beta(B/p)\alpha=-c$, and for distinct singletons it is
$-\beta^2\alpha=-c$. All diagonal correlations are one.
Common multiplication by $\omega$ preserves the joint law and
all block products. Since $p$ is odd, $\omega^2\ne1$; therefore
every holomorphic second moment is zero. We have obtained
\eqref{eq:ternary-odd-covariance} for the full odd-mode system
at a value $c>c_0$.

If $r=1$, then $s=3$, so necessarily $p=3$ and $k=2$.
This is exactly the separately treated $P_8$ case above.
For all other parameters the preceding construction applies.

Finally mix the constructed measure with Haar measure in proportion
$c_0/c$ to obtain the exact target $c_0$ and a positive Haar
component. Distinct odd modes have neither equal nor opposite
frequencies, so the relevant Haar moments vanish. The smoothing,
finite-character correction and finite-time Gram argument above
apply on the now identified phase torus without change. They produce
an actual bounded compactly supported time density, and the
spectral covariance identity gives $J/(p^k-1)$.
\end{proof}

The new cases include $P_8,P_{24},P_{26},P_{48},P_{80},P_{120},
P_{124},\ldots$. Prime moduli themselves were already covered by
\cite{Kang2026}; the added assertion is the entire family of higher
odd prime powers. Neither this theorem nor the preceding ternary
specialization settles all remaining composite moduli. The bounded
priority statement above applies to this stronger family as well.
\section{Reproducible integer certificates}\label{app:code}
The following Python programs are embedded in this single source file.
They can be copied from the \texttt{lstlisting} environments and run
independently. The main theorem requires no external data file or
computer search. The first verifier uses only the Python standard library.
Its default run includes the compact certificate for $m=127$.

\subsection{Full cycle and digit-tree verification}
The depth chosen by this implementation is conservative. A lower depth
satisfying the displayed size inequalities is equally valid.

\begin{lstlisting}[language=Python]
from math import ceil, log2

def v2(n):
    assert n
    return (abs(n)&-abs(n)).bit_length()-1

def weight(n):
    return 1 << v2(n)

def certificate(p):
    assert p>1 and p%2
    L=(p-1).bit_length()
    s=v2(p*p-1)
    r=v2(p-1)
    t=1<<(L+3)
    J=L+s+2
    q=1<<J
    a0=1<<r
    D=t*(p-1-a0)
    exc=[]
    if p%4==1:
        for M in range(D.bit_length()):
            if D>>M&1:
                assert M>s
                exc.append(p+(1<<M))
    elif p==3:
        assert D==0
    else:
        C=(1<<s)-4
        k=1<<(L+1)
        B=D-k*C
        assert B>0 and B%(1<<(s+1))==0
        for M in range(B.bit_length()):
            if B>>M&1:
                assert M>s
                exc.append(p+(1<<M))
        assert len(exc)<=k
        exc += [-1]*(k-len(exc))
    assert 2*len(exc)<=t
    h=[1]*t
    for i,A in enumerate(exc): h[2*i+1]=A
    g=[h[(i+1)%t]-p*h[i] for i in range(t)]
    assert all(x and x%2 for x in h)
    assert all(g)
    assert sum(map(weight,g))==(p-1)*sum(map(weight,h))
    # Build all cycles modulo 2^J, paired by negation.
    H={}
    for u in range(1,q,2):
        if u in H:continue
        x=u
        for i in range(t):
            assert x not in H and (-x)%q not in H
            H[x]=q*h[i]
            H[(-x)%q]=-q*h[i]
            x=p*x%q
        assert x==u
    G={u:H[p*u%q]-p*H[u] for u in H}
    assert len(H)==q//2
    assert all(H[q-u]==-v for u,v in H.items())
    assert sum(map(weight,G.values()))==(p-1)*sum(map(weight,H.values()))
    a=max(J,(4*max(max(map(abs,H.values())),max(map(abs,G.values())))+1).bit_length())
    # This a is conservative: condition is 4*max+1 < 2^(a+1).
    assert 4*max(map(abs,H.values()))+1<1<<(a+1)
    assert 4*max(map(abs,G.values()))+1<1<<(a+1)
    assert a<=5*L+10
    X={u:1+4*v for u,v in H.items()}
    Y={u:-1+4*v for u,v in G.items()}
    assert all(Y[u]+p*X[u]-X[p*u%q]==p-2 for u in H)
    # Unique signed binary digits for odd |x|<2^(a+1).
    def digits(x):
        z=(x+(1<<(a+1))-1)//2
        return [2*((z>>i)&1)-1 for i in range(a+1)]
    def layer_sums(X):
        sums=[]
        for i in range(a+1):
            b=max(1,min(i,J))
            d={}
            for u,x in X.items():
                k=u% (1<<b)
                z=digits(x)[i]
                assert k not in d or d[k]==z
                d[k]=z
            multiplicity=1 if i<=J else 1<<(i-J)
            sums.append(multiplicity*sum(d.values()))
        return sums
    SH=layer_sums(X);SG=layer_sums(Y)
    assert sum(SH)==sum(SG)==0
    # e_0=a, e_i=a-i-1 for i>=1.
    e=[a]+[a-i-1 for i in range(1,a+1)]
    assert sum(ei*(gi+(p-1)*hi) for ei,gi,hi in zip(e,SG,SH))==0
    return dict(p=p,L=L,s=s,J=J,t=t,excursions=len(exc),a=a,N=f'{p}*2^{a}',
                maxH=max(map(abs,H.values())),maxG=max(map(abs,G.values())))

if __name__=='__main__':
    import sys,json
    for p in map(int,sys.argv[1:] or [3,5,7,11,13,17,31,127]):
        print(json.dumps(certificate(p)),flush=True)

\end{lstlisting}

\subsection{Independent arithmetic audit}
This implementation uses a recursive signed-digit decoder and sums
excursion increments without constructing the full digit trees.
Its final optional numerical block checks the scalar in the elementary
transfer lemma with \texttt{mpmath}; these numerical checks are not used
as a proof of vanishing. All preceding checks use exact integers.

\begin{lstlisting}[language=Python]
from fractions import Fraction
from math import log2
import mpmath as mp

def v(n):
    assert n
    n=abs(n); r=0
    while n%2==0:n//=2;r+=1
    return r

def digits(n,a):
    # Recover recursively using lowest signed digit congruence, independently.
    out=[]
    for i in range(a):
        d=1 if n%4==3 else -1
        out.append(d);n=(n-d)//2
    assert n in (-1,1)
    return out+[n]

def lamnum(n,a):
    d=digits(n,a)
    return 2*a*d[0]+sum((a-i-1)*(1<<i)*d[i] for i in range(1,a+1))

pair_count=0
for a in range(2,12):
    for h in range(1,1<<(a-1)):
        if 4*h+1>=1<<(a+1): continue
        for e in [-1,1]:
            assert lamnum(e+4*h,a)+lamnum(e-4*h,a)==-e*(1<<(v(h)+3))
            pair_count+=1
print('pairing exact tests',pair_count)

for m in range(3,10000,2):
    L=(m-1).bit_length();s=v(m*m-1);r=v(m-1)
    t=1<<(L+3);J=L+s+2;b=1<<r
    if m%4==1:
        D=t*(m-1-b)
        E=[m+(1<<j) for j in range(D.bit_length()) if D&(1<<j)]
    elif m==3:E=[]
    else:
        k=1<<(L+1);C=(1<<s)-4;D=t*(m-3);B=D-k*C
        assert B==t*(m-2-(1<<(s-2)))>0
        E=[m+(1<<j) for j in range(B.bit_length()) if B&(1<<j)]
        E.extend([-1]*(k-len(E)))
    assert 2*len(E)<=t
    total=t*b
    maxh=max([1]+[abs(x) for x in E]);maxg=m-1
    for x in E:
        total+=(1<<v(x-m))+(1<<v(1-m*x))-2*b
        maxg=max(maxg,abs(x-m),abs(1-m*x))
    assert total==(m-1)*t
    a=J+3*L+7
    assert a<=5*L+10
    assert 4*(1<<J)*max(maxh,maxg)+1<1<<(a+1)
    assert pow(m,t,1<<J)==1 and pow(m,t//2,1<<J)!=1
    # A subgroup containing -1 would have its unique order-two element equal -1.
    assert pow(m,t//2,1<<J)!=(1<<J)-1
print('cycle and bound audits: all odd m 3..9999')

# Independent exact elementary logarithm transfer, numerically cross-checking only constants.
mp.mp.dps=70
for a in range(2,8):
    q=1<<a
    levels=[[3]]+[[2*i-3+(u%3) for u in range(1,1<<i,2)] for i in range(1,a)]
    K=-3
    top=[]
    for u in range(1,q,2):
        top.append(Fraction(K-levels[0][0]-sum((1<<i)*levels[i][((u%(1<<i))-1)//2] for i in range(1,a)),q))
    levels.append(top)
    vals=[]
    for n in range(1,q+1):
        if n==q:vals.append(Fraction(levels[0][0]));continue
        i=a-v(n);u=n//(1<<v(n));vals.append(Fraction(levels[i][(u-1)//2]))
    assert sum(vals)==0
    S=list(map(sum,levels))
    B=a*S[0]+sum((a-i-1)*S[i] for i in range(1,a+1))
    L=-sum(mp.mpf(c.numerator)/c.denominator*mp.digamma(mp.mpf(n)/q) for n,c in enumerate(vals,1))/q
    predicted=-mp.log(2)*mp.mpf(B.numerator)/B.denominator/q
    assert abs(L-predicted)<mp.mpf('1e-60')
print('digamma transfer cross-check: a=2..7')

\end{lstlisting}

\subsection{Graph deletion checks and the source-budget table}
The next program is the exact labelled-core check described in the graph
appendix. It recomputes graph indices using rational numbers. A full run
through core order six may take appreciably longer than the arithmetic
checks above.

\begin{lstlisting}[language=Python]
from fractions import Fraction as Q
from itertools import combinations
from math import factorial
import json

def m(adj):
    d=list(map(len,adj))
    return sum((Q(1,max(d[u],d[v])) for u in range(len(adj)) for v in adj[u] if u<v),Q(0))
def drop_vertices(adj,remove):
    return [set(x)-remove if i not in remove else set() for i,x in enumerate(adj)]
def check(adj,n):
    d=list(map(len,adj)); coreedges=[(u,v) for u in range(n) for v in adj[u] if u<v and v<n]
    val=m(adj)
    for u,v in coreedges:
        a=[set(x) for x in adj];a[u].remove(v);a[v].remove(u)
        if m(a)<=val:return False
    if min(d[:n])<3:return False
    src=[u for u in range(n) if all(d[v]>=d[u] for v in adj[u] if v<n)]
    leaves={next(v for v in adj[u] if v>=n) for u in src}
    assert all(not(set(src)&adj[u]) for u in src)
    budget=sum((Q(1,d[u]) for u in src),Q(0))
    assert val-m(drop_vertices(adj,leaves))==budget
    full=val-m(drop_vertices(adj,set(range(n,len(adj)))))
    assert full>=budget
    D=max(d);F={D:1}
    for j in range(D-1,2,-1):F[j]=1+(j-1)*F[j+1]
    assert n<=sum(F[d[u]] for u in src)
    if full<1:assert n<=3*sum(factorial(j) for j in range(1,D-1))-1
    return True

counts={}; tested=0
for n in range(3,7):
    poss=list(combinations(range(n),2));passed=0;decorated=0
    for mask in range(1<<len(poss)):
        a=[set() for _ in range(n)]
        for i,(u,v) in enumerate(poss):
            if (mask>>i)&1:a[u].add(v);a[v].add(u)
        if min(map(len,a))<2:continue
        for lmask in range(1<<n):
            b=[set(x) for x in a]
            for u in range(n):
                if (lmask>>u)&1:
                    b[u].add(len(b));b.append({u})
            decorated+=1
            if check(b,n):passed+=1
    tested+=decorated;counts[n]={'decorations':decorated,'all_supportive':passed}
print(json.dumps({'tested':tested,'by_core_order':counts,'status':'all exact checks passed'},indent=2))

\end{lstlisting}

The table of integer improvements is reproduced by this short program.
\begin{lstlisting}[language=Python]
from math import lcm, factorial
for D in range(4,9):
    F={D:1}
    for d in range(D-1,2,-1):
        F[d]=1+(d-1)*F[d+1]
    L=lcm(*range(3,D+1))
    dp=[-1]*L; dp[0]=0
    for w in range(L):
        if dp[w]<0: continue
        for d in range(3,D+1):
            v=w+L//d
            if v<L: dp[v]=max(dp[v],dp[w]+F[d])
    print(D,3*sum(factorial(j) for j in range(1,D-1))-1,max(dp))
\end{lstlisting}

\section{Source and dependency record}
The positive arithmetic construction uses only the explicit lemmas proved
in Sections 2--6. The pure-power exclusion uses the published Okada unit
equations, and its reduction to the alternating harmonic series is proved
here. No theorem about Euler's constant, no assumed independence of
Dirichlet $L$-values, and no unverified numerical zero enters that argument.
The graph order theorem is independent of the corpus; only its radius
corollary invokes the earlier structure reduction. The Steklov argument
uses standard variational and Sobolev facts whose exact applications
are given in its proof.

The quantum-walk proof identifies the full integer frequency-relation
lattice and constructs a measure on its constrained torus. Its final
finite-time correction is included explicitly, so a limiting torus law
is not being substituted for a distribution on actual observation times.

The initial research screen also rejected proposed new targets when an
identical or broader result was already present in recent primary
literature. Those abandoned directions are not counted as results of
this manuscript. The bibliography below lists the sources actually used
for the stated theorems and their provenance.

\begin{thebibliography}{99}\small
\bibitem{MP2022}
M.~Ram Murty and S.~Pathak,
\emph{The Okada space and vanishing of $L(1,f)$},
Functiones et Approximatio Commentarii Mathematici
\textbf{66} (2022), no.~1, 35--57,
\href{https://doi.org/10.7169/facm/1952}{doi:10.7169/facm/1952}.
Inspected author manuscript dated 7 April 2021:
\url{https://www.cmi.ac.in/~siddhi/Okada_space_RM_SP_Final.pdf}.
Theorem 1.1 gives the necessary unit equations; Section 4.4 states the
period-classification question and Theorem 4.8.

\bibitem{Okada1982}
T.~Okada, \emph{On a certain infinite series for a periodic arithmetical
function}, Acta Arithmetica \textbf{40} (1982), 143--153.
The needed criterion is used through its explicit restatement in
\cite[Theorem 1.1]{MP2022}.

\bibitem{Corpus093}
Research audit XXIV, corpus item [093],
\emph{When can $\sum f(n)/n$ vanish for a periodic $f$ with values $\pm1$?},
file \texttt{[093] pm1\_vanishing\_paper.tex}, research draft.
Theorem E and Conjecture 9.1, in the repository snapshot specified
in Section 1.
\url{https://github.com/QEEDDD/Verifiable-Mathematical-Reasoning-Corpus}.

\bibitem{Corpus087}
VMRC, corpus item [087],
\emph{Spanning trees of small diameter and the Randi\'c index},
file \texttt{[087] fajtlowicz\_radius\_core\_reduction.tex}, research draft.
Proposition \texttt{prop:structure} and Remark \texttt{rem:D6};
same repository snapshot as \cite{Corpus093}.

\bibitem{Corpus094}
VMRC, corpus item [094], \emph{A bag reduction for the Randi\'c--radius problem},
file \texttt{[094] Randic\_radius\_bag\_reduction.tex}, research draft;
same repository snapshot as \cite{Corpus093}.

\bibitem{CPS}
M.~Cygan, M.~Pilipczuk, and R.~\v{S}krekovski,
\emph{On the inequality between radius and Randi\'c index for graphs},
MATCH Communications in Mathematical and in Computer Chemistry
\textbf{67} (2012), 451--466.
\url{https://match.pmf.kg.ac.rs/electronic_versions/Match67/n2/match67n2_451-466.pdf}.

\bibitem{Lagace2020}
\emph{Open problems from the miniconference on sharp eigenvalue estimates
for partial differential operators}, April 2020,
Open Problem 10, proposed by Jean Lagac\'e, printed page 3.
\url{https://publish.illinois.edu/eigenvalues2020/files/2020/04/Open-problems.pdf}.
The complete four-page problem list was inspected.

\bibitem{GP2017}
A.~Girouard and I.~Polterovich,
\emph{Spectral geometry of the Steklov problem},
Journal of Spectral Theory \textbf{7} (2017), 321--359.
\href{https://doi.org/10.4171/JST/164}{doi:10.4171/JST/164}.
Example 4.2.5 treats annuli by separation of variables.
\url{https://ems.press/content/serial-article-files/33609}.
\bibitem{CoutinhoGuo2024}
G.~Coutinho and K.~Guo,
\emph{Selected open problems in continuous-time quantum walks},
Special Matrices \textbf{12} (2024), 20240025.
\href{https://doi.org/10.1515/spma-2024-0025}{doi:10.1515/spma-2024-0025}.
Section 4.2 discusses randomized observation times.

\bibitem{Kang2026}
M.~Kang, \emph{Arithmetic of Bohr Frequencies Governs Uniform Mixing
in Randomly Timed Quantum Spin Chains}, arXiv:2609.14463v1,
13 September 2026.
\url{https://arxiv.org/abs/2609.14463v1}.
The inspected version establishes several positive and negative families
and leaves other moduli outside its classification.
\end{thebibliography}
\end{document}
